\documentclass[12pt,a4paper]{article}
\usepackage[romanian]{babel}
\usepackage{fontspec}
\setmainfont{Latin Modern Roman}
\usepackage{geometry}
\geometry{margin=2.5cm}
\usepackage{graphicx}
\usepackage{enumitem}
\usepackage{longtable}
\usepackage{framed}

\title{AMSS 2026 — Examen restanță (model) \\ \large Soluții și barem orientativ}
\author{Traian-Florin Șerbănuță}
\date{2026}

\begin{document}
\maketitle

\begin{framed}
\noindent Acestea sunt \textbf{soluții model pentru subiectul-model}
(domeniul bike-sharing din `examen-2026.pdf`). Subiectul real de la
restanță folosește \textbf{alt scenariu} — dar \textbf{baremul și modul
de notare sunt aceleași}. Folosește acest document ca să înțelegi
\emph{ce se așteaptă} de la tine, nu ca listă de răspunsuri de memorat.
\end{framed}

\bigskip

\section*{Cum se notează (barem orientativ)}

Examenul oglindește rubrica apărării orale (F1 + F3 + F4). Se notează
\textbf{competența}, nu prezentarea. Pondere orientativă:

\begin{description}[leftmargin=4.5em,style=nextline]
  \item[Partea I (F1) — \char`~50\%] Identificarea defectelor: cel puțin
        cinci defecte reale, fiecare cu \emph{categorie}, \emph{severitate}
        și \emph{corecție}. Răspuns slab: „arată ciudat" fără categorie
        sau corecție.
  \item[Partea II (F3) — \char`~25\%] O \emph{regulă de domeniu concretă}
        plus un \emph{criteriu clar de respingere}. Răspuns slab: „aș cere
        AI-ului o diagramă corectă", fără regulă și fără criteriu.
  \item[Partea III (F4) — \char`~25\%] Identificarea inconsistenței
        \emph{plus} parcurgerea lanțului de trasabilitate \emph{plus}
        corecția. Răspuns slab: „multiplicitatea e greșită", fără traseu.
\end{description}

\section*{Partea I — Critică (F1): soluție model}

Defectele plantate în artefactele-model (minim cinci așteptate):

\paragraph{Diagrama de clase.}

\begin{longtable}{p{4.2cm} p{1.8cm} p{6cm}}
\textbf{Defect (categorie)} & \textbf{Severitate} & \textbf{Corecție} \\
\hline
\endhead
clasă-zeu: \texttt{BikeSystem} (ține tot + \texttt{doEverything()}) &
critică & distribuie responsabilitățile pe \texttt{Rental},
\texttt{Bike}, \texttt{Station}; elimină clasa-zeu \\
multiplicitate greșită: \texttt{Rental "*" -- "*" Bike} & critică &
R1: o închiriere = o bicicletă → \texttt{Rental "*" -- "1" Bike} \\
agregare lipsă: \texttt{Station -- Bike} (linie simplă) & majoră &
R2: stația \emph{deține} biciclete care îi supraviețuiesc →
\texttt{Station o-- Bike} \\
asociere vagă: \texttt{Rental -> Payment} fără multiplicitate & minoră &
precizează multiplicitatea (o închiriere are o plată) \\
\end{longtable}

\paragraph{Diagrama de stare (Bike).}

\begin{longtable}{p{4.2cm} p{1.8cm} p{6cm}}
\textbf{Defect (categorie)} & \textbf{Severitate} & \textbf{Corecție} \\
\hline
\endhead
stare orfană / cerință nemodelată: nimic nu intră în
\texttt{Maintenance} (R3 cere ca o bicicletă cu defect să intre în
mentenanță) & majoră & adaugă \texttt{InUse --> Maintenance : fault} \\
lipsă pseudo-stare inițială: nu există \texttt{[*] --> Available} &
majoră & adaugă starea inițială \\
\end{longtable}

\section*{Partea II — Raționament (F3): răspuns model}

\textit{R3: o bicicletă cu defect intră în mentenanță și revine
disponibilă după reparație.}

\begin{quote}
Aș specifica regula de domeniu explicit: „din \texttt{InUse}, o
bicicletă cu defect trece în \texttt{Maintenance}; din
\texttt{Maintenance} revine în \texttt{Available} după reparație;
mașina are o stare inițială (\texttt{Available})." Aș cere ca fiecare
stare să fie \emph{accesibilă} și să aibă o ieșire.

Aș \textbf{respinge}: o stare \texttt{Maintenance} fără tranziție de
intrare (orfană); absența stării inițiale; tranziții fără eveniment;
stări inventate care nu derivă din R3.
\end{quote}

Se punctează: regula de domeniu concretă \emph{și} criteriul de
respingere. „Aș cere o diagramă bună" nu primește punctaj.

\section*{Partea III — Trasabilitate (F4): răspuns model}

\textit{R1: o închiriere este pentru exact o bicicletă.}

\begin{quote}
R1 spune \emph{exact o bicicletă}. Diagrama arată
\texttt{Rental "*" -- "*" Bike} (mulți-la-mulți): o închiriere ar putea
referi mai multe biciclete. \textbf{Inconsistență directă.}

Lanțul de trasabilitate: R1 (cerință) → use case „Închiriază
bicicletă" → clasa \texttt{Rental} și relația ei cu \texttt{Bike}.
Citind multiplicitatea cu voce tare — „o închiriere are mai multe
biciclete?" — contrazice R1. Corecție: \texttt{Rental "*" -- "1" Bike}.
\end{quote}

Se punctează: identificarea inconsistenței \emph{plus} parcurgerea
lanțului cerință→model \emph{plus} corecția.

\end{document}
